% Adapted from the supplied PA194N no-push proof, with formula coefficient six.
% This is an ordinary mathematical proof, not a new proof-assistant endpoint.
\section{The radix-six no-push command compiler}
\label{sec:np-controller}\label{sec:commands}

\subsection{Exact numerical operations and command contracts}

Let
\[
 p(a,b)=\frac{(a+b)(a+b+1)}2+b,\qquad [a;w]=p(a,w),\qquad
 I(a,b)=6p(a,b),\qquad U(i,a)=6p(a,i)+2.
\]
Iterated brackets put their leftmost entry at the top. The Cantor
projections are total and satisfy
\[
 \operatorname{hd}[a;w]=a,\quad \operatorname{tl}[a;w]=w,
 \quad[0;0]=0,\quad\operatorname{hd}(0)=\operatorname{tl}(0)=0.
\]
A numerical stack is therefore not being treated as a free list.
A command store is $(t,w,d,h)$: current formula, stack, offset, and cached
formula. Code zero is truth and code one is the target contradiction.
The new MP command, with $y=d+x$, is
\begin{equation}\label{eq:np-mp}
 \mathsf{MP}_x(t,w,d,h)=
 \begin{cases}
 (y,w,y,0),& I(t,y)=h,\\
 (0,w,y,0),& I(t,y)\ne h.
 \end{cases}
\end{equation}
Shift is retained literally:
\begin{equation}\label{eq:np-shift}
 S(t,w,d,h)=
 \begin{cases}
 \mathsf{halt},&t+h=1,\\
 (h,\operatorname{tl}(w),d,\operatorname{hd}(w)),&t+h\ne1.
 \end{cases}
\end{equation}
For $D=d+x$, $H=h+y$, the unchanged AX command has output
\begin{equation}\label{eq:np-ax}
 \mathsf{AX}_{A,x,y}(t,w,d,h)
   =\bigl(g(H,D,A),[q_1,\ldots,q_9,t;w],0,0\bigr).
\end{equation}
Here $q_1=\Gamma(A)=I(A,U(A,A))$, $q_2,\ldots,q_7$ are the
six retained head schemas in reverse emission order, interpreted with
the radix-six constructors defined in Section~\ref{sec:calculus},
\[
 q_8=\rho(D,A,H,d)=I(I(I(H,d),D),I(I(D,H),I(A,H))),
 \qquad q_9=U(d,h).
\]
The gate $g$ is the separate-head guarded separation gate at coefficient six.
The inherited free-constructor emitter audit is retained in
\texttt{construction/inputs/parent6.emitter\_audit.json}.
The new source splice occurs wholly after that emitter and retains all
of its resolved instructions and successor addresses under relocation. All $q_1,\ldots,q_8,g$ are divisible by six,
and $q_9\equiv2\pmod6$, with $q_9\ge2$. In particular none is code one.
This is a statement about constructor values, requiring no consistency
assumption and holding for arbitrary incoming natural stores.

The count in (\ref{eq:np-ax}) is nine generated stack entries plus the
saved old top. The two identities needed below are
\begin{align}
 \mathsf{AX}_{A,x,y};S^9(t,w,d,h)
   &=(q_8,[t;w],0,U(d,h)),\label{eq:np-nine}\\
 \mathsf{AX}_{A,x,y};S^{10}(t,w,d,h)
   &=(U(d,h),w,0,t).\label{eq:np-ten}
\end{align}
All displayed Shifts are nonhalting: their tests involve two generated
non-one codes, zero and a generated code, or the universal code of value
at least two. This also establishes exact stack effects when zero pairs
collapse.

\subsection{Uniform stack operations and a one-command discard}

Write $A_0=\mathsf{AX}_{0,0,0}$ and define
\[
 \mathsf{Push}=A_0;S^9;\mathsf{MP}_0,\qquad
 \mathsf{Pop}=A_0;S^{11},\qquad
 \mathsf{Discard}=\mathsf{MP}_0,\qquad
 \mathsf{Offset}_i=\mathsf{MP}_i.
\]
Their exact contracts are
\begin{align}
 \mathsf{Push}(t,w,d,h)&=(0,[t;w],0,0),\label{eq:np-push}\\
 \mathsf{Pop}(t,[a;w],d,h)&=(t,w,0,a),\label{eq:np-pop}\\
 \mathsf{Discard}(t,w,0,h)&=(0,w,0,0),\label{eq:np-discard}\\
 \mathsf{Offset}_i(0,w,0,0)&=(0,w,i,0).\label{eq:np-offset}
\end{align}
For Push, (\ref{eq:np-nine}) makes the final MP compare an implication
with a universal code, so it fails with offset zero. For Pop, the first
ten Shifts give (\ref{eq:np-ten}). The eleventh Shift tests
$U(d,h)+t\ge2$, hence it does not halt even when $t=1$.
It pops exactly the displayed item. Thus Pop is total on all natural
stores. In particular $\mathsf{Pop}(t,0,d,h)=(t,0,0,0)$ and
$\mathsf{Push}(0,0,d,h)=(0,0,0,0)$.

For Discard, the two MP outcomes in (\ref{eq:np-mp}) coincide when
$d=x=0$, for arbitrary $h$: both are $(0,w,0,0)$. This gives a single-command discard,
where the old dummy-push controller needed a longer word. The old cache
need not be zero for this identity. For Offset,
if $i>0$ then $I(0,i)>0$, so the comparison fails; if $i=0$ it succeeds
and returns $(0,w,0,0)$. The same one-command definition covers both cases.
These simplifications affect the proof-word compiler; they do not claim
additional physical state savings.

\subsection{Axiom selectors and generalization}

In the following contracts the arbitrary incoming store is $(t,w,0,0)$.
For $1\le j\le7$, the word
\[
 \mathsf{Discard};\mathsf{AX}_{A,D,H};S^{j+1};\mathsf{Pop}^{9-j}
\]
returns $(q_j,w,0,0)$ with the parameters evaluated at the desired
$A,D,H$. After the initial Shifts, the remaining $9-j$ stack entries
are exactly the unselected generated tail and the saved zero. The Pops
preserve the selected top and consume those entries; the last pop puts
that saved zero into the cache. For the dominated head, if $u\le a$,
use
\[
 \mathsf{Discard};\mathsf{Offset}_u;
 \mathsf{AX}_{b,a-u,c};S^9;\mathsf{Pop}.
\]
It returns $(\rho(a,b,c,u),w,0,0)$, with the numerical dominance condition
coming from $a=u+(a-u)$. A permitted gate is obtained by
\[
 \mathsf{Discard};\mathsf{AX}_{A,D,H};\mathsf{Pop}^{10},
\]
which returns $(g(H,D,A),w,0,0)$. Truth uses Discard alone.

The word
\[
 \mathsf{Gen}_i=\mathsf{Push};\mathsf{Offset}_i;S;
                \mathsf{AX}_{i,0,0};S^{10}
\]
returns $(U(i,a),w,0,0)$ from $(a,w,0,0)$. Before its AX the store is
$(0,w,i,a)$; its last ten Shifts use (\ref{eq:np-ten}) with saved top zero.
Its first Shift tests $0+0$, and the later ones are among the safe AX
shifts. Hence this generalization word is total even for $a=1$.

\subsection{Modus ponens with its spare zero supplied by the compiler}

Suppose recursively compiled proofs return $a$ and $I(a,b)$ with the
stack framed and both caches cleared, or halt earlier. The MP compiler is
\begin{equation}\label{eq:np-compile}
 \mathsf{Discard};\mathsf{Push};C(I(a,b));\mathsf{Push};
 C(a);\mathsf{Push};S;S;\mathsf{MP}_b;\mathsf{Pop}.
\end{equation}
The first two commands reserve a zero below the premise computations.
If the recursive calls have returned, the store before the two Shifts is
\[
             (0,[a,I(a,b),0;w],0,0).
\]
The first Shift returns $(0,[I(a,b),0;w],0,a)$.
The second halts if $a=1$; otherwise it returns
$(a,[0;w],0,I(a,b))$. MP succeeds and returns
\[
                         (b,[0;w],b,0).
\]
By the total Pop contract, the final word returns $(b,w,0,0)$, including
$b=1$. The stack zero is now consumed once, at the compiler level.
At $a=b=w=0$ the same proof uses the literal Cantor projection identities;
no extra distinction between a zero push and an empty stack is introduced.

Induction on a finite derivation $p$ in the inherited separate-head
calculus therefore proves:
\begin{quote}
From every $(t,w,0,0)$, $C(p)$ either halts or returns
$(\operatorname{concl}(p),w,0,0)$.
\end{quote}
The axiom and generalization cases use the words above. The MP case is
(\ref{eq:np-compile}); extending a word after an early halt still halts.
In particular a refutation compiler ends in top one, if it has not already
halted. Its final appended Shift halts because $1+0=1$.
For arbitrary incoming $t,w,d,h$, first apply $A_0$ to clear $d,h$.
Thus the concrete word
\[
                         A_0;C(p);S
\]
halts from \emph{every} incoming store whenever $p$ derives code one.
The compiler is a primitive-recursive traversal of the finite derivation,
and PA proves its correctness by induction on that coded derivation.
It does not assume a reset between the scheduler's unsuccessful trials.

The unchanged scheduler enumerates finite command words with the usual
Cantor continuation coding. A cofinal family of halting fresh candidates
is obtained by prefixing the universal word by
$\mathsf{AX}_{k,0,0}$. Let $E(\epsilon)=0$ and let a positive continuation be encoded by
$E(c::w)=1+\operatorname{encodeCommand}(c,E(w))$. In particular
$\operatorname{encodeCommand}(\mathsf{AX}_{A,x,y},z)
=p(A+2,p(y,p(x,z)))$. Put $z=E(A_0;C(p);S)$, using this positive
continuation convention rather than a fresh-candidate code. The fresh
candidate codes for the prefixed words are exactly
$p(k+2,p(0,p(0,z)))$, which are strictly increasing and
unbounded in $k$. The fresh counter begins at two, so omitting smaller
candidates loses none of this cofinal family. Every decoded continuation
strictly decreases, and the unchanged source macro loops terminate.
The resulting finite-word and finite-trace construction is formalizable
in PA by primitive recursion and induction.

\subsection{Soundness and the PA calibration}

Annotate the current top $t$, cached head $h$, and every live canonical
Cantor-stack entry by a derivation in the separate-head calculus, or by a
PA proof of its arithmetic translation. For $w>0$, its decoded tail is
strictly smaller than $w$; thus the canonical stack and its annotation
have a primitive-recursive bounded construction. A zero stack supplies
truth when popped, and pushing zero onto zero preserves the empty
canonical annotation.

AX adds the unchanged valid axiom instances, a generalization of the old
cached theorem, and the saved old top. On a successful MP, the annotations
of $t$ and $I(t,y)=h$ produce an annotation for $y$. The stack is retained,
and the cache receives the truth proof. Failure outputs truth and retains
the stack. Shift transports the annotations. If its test succeeds, then
$t+h=1$ implies that one of the two annotated formulas is code one.
The inherited syntactic arithmetic translation turns that annotation into
a PA refutation. No annotation is imposed on the arithmetic offset $d$.

Consequently PA proves both finite syntactic reductions:
\[
 \text{source halt}\ \longrightarrow\ \text{PA refutation},
 \qquad
 \text{PA refutation}\ \longrightarrow\ \text{source halt}.
\]
For the second reduction, first use the supplied PA-to-radix-four
translation and then the hereditary recoding and restricted-axiom proof
templates of Section~\ref{sec:calculus} to obtain a radix-six
refutation. Apply the all-store compiler above to that refutation.
For the first, compose the proof annotation with the radix-six-to-PA
translation of that same section. This uses the existing logical calibration,
including its treatment of the code-zero alias and coincident variable
labels. Nothing in the no-push modification changes that calculus or
requires a new logical axiom, a reflection principle, or consistency as
an additional assumption.

